\chapter{Conclusion}

This survey set out to do one thing, to lay the esoteric models of reality side by side and see what they
held. Two findings stand out.

The first is that the models are real models. They are not vague feelings dressed in story. Each names a
source, a process of unfolding, a structure of worlds, an account of the human, and a path. The two studied
in most depth, Kabbalah and Tibetan Buddhism, are as detailed and internally consistent as any formal system,
and the others are far richer than a casual reading suggests. Taken together they form an inherited atlas of
the inner and cosmic order, built over thousands of years by people who gave their lives to the work.

The second is that the models converge. Beneath very different images and vocabularies, the same sequence
recurs. There is one source. It goes out into the many. The many lose sight of it. And the work of a life,
and the point of a path, is to find the way back, to remember, to be gathered home, to tune again to the one.
The traditions disagree on details, and a few disagreements are worth keeping as live alternatives. But the
shared shape is stable across cultures that had no contact with one another, which is the strongest reason to
think they were tracking something real.

That shared something is the reason this survey was gathered. To build a serious account of consciousness,
the inner life, the sense of the divine, and the worlds of beings both bright and dark, it helps to know the
full depth and scope of what was already mapped, so that nothing important is missed. The traditions have
already charted the source, the descent, the forgetting, the human turn, and the return. The next work is to
take that inherited map seriously, to test it, and to follow the path it describes, the return to the source
of reality, the flow behind the field, the way that lies beyond the senses.
